\subsection{Symbolic Mapping Invariants}
\label{app:symbolic-mapping-reassociation}
The mapping state is $\Sigma=(V,\mathcal M,\Pi)$, where
$V[t]=(\operatorname{Ord}_t,e_t)$ stores an ordered interface and a value expression.
A mapping $(\operatorname{Ord}_t,\vec G,q)\in\mathcal M$ satisfies
\begin{equation}
 \gath(\operatorname{Ord}_t,q)=\vec G.
 \label{extended:eq:mapping-invariant}
\end{equation}
Each partition in $\Pi$ lists pairs $(\vec G_i,q_i)$ that cover an
interface exactly once: the concatenated position lists are a permutation
of its positions. Different partitions may group the same interface
differently without changing the tensor. A producer-first traversal
aligns operand interfaces, transports reads through validated mappings,
and composes the declared output axes and value semantics.

\paragraph{Reconstruction and composition.}
For a complete binding, concatenate its groups and their position lists.
If $q=q_1\Vert\cdots\Vert q_\ell$ is a permutation, then
\begin{equation}
 Q=\gath(\vec G_1\Vert\cdots\Vert\vec G_\ell,\operatorname{inverse}(q))
 \label{extended:eq:rebuilt-interfaces}
\end{equation}
restores tensor dimension order, including interleaved groups.
Paired overlap matching refines partitions; gather composition
(Equation~\eqref{eq:gather-composition}) derives further selections.
Lemma~\ref{lem:validated-mappings} proves preservation of these invariants.

\paragraph{From mappings to attribute programs.}
\label{app:mappings-to-attribute-programs}
For admitted plans, the symbolic adapter interprets the shared metadata IR
of Theorem~\ref{thm:symbolic-compilation} through these mappings. Source-read
correspondence supplies logical axes; restricted additive components are
matched independently. Unresolved reductions or ambiguous read assignments
make lowering unsuccessful, leaving the residual program available for
native interpretation. Supplied extent and coordinate expressions are
preserved, and completed candidates require independent validation.
Appendix~\ref{app:batched-reassociation-program} gives the complete
reassociation case study. Traversal pseudocode and expanded tutorials are
archived in \nolinkurl{evaluation/supplementary-materials/}.
